\documentclass[conference]{IEEEtran}
\IEEEoverridecommandlockouts

\usepackage{cite}
\usepackage{amsmath,amssymb,amsfonts}
\usepackage{algorithmic}
\usepackage{graphicx}
\usepackage{textcomp}
\usepackage{xcolor}
\usepackage{booktabs}
\usepackage{array}
\usepackage{url}
\usepackage{hyperref}
\usepackage{longtable}
\usepackage{tabularx}

\def\BibTeX{{\rm B\kern-.05em{\sc i\kern-.025em b}\kern-.08em
    T\kern-.1667em\lower.7ex\hbox{E}\kern-.125emX}}

\begin{document}

\title{Proposal for CrisisEcho: LLM-Augmented Early Detection\\
of Local Emergencies from Social Media}

\author{
\IEEEauthorblockN{Samuel Enam Zih}
\IEEEauthorblockA{
Computer Science\\
Virginia Tech\\
Alexandria, VA, USA\\
\texttt{szih@vt.edu}\\
(LEADER)}
\and
\IEEEauthorblockN{Christopher Drake Williams}
\IEEEauthorblockA{
Computer Science\\
Virginia Tech\\
Alexandria, VA, USA\\
\texttt{cdwilliams@vt.edu}}
\and
\IEEEauthorblockN{Arjun Sivakumar}
\IEEEauthorblockA{
Computer Science\\
Virginia Tech\\
Alexandria, VA, USA\\
\texttt{arjuns@vt.edu}}
\and
\IEEEauthorblockN{Ki Hong Park}
\IEEEauthorblockA{
Computer Science\\
Virginia Tech\\
Alexandria, VA, USA\\
\texttt{kihong@vt.edu}}
\and
\IEEEauthorblockN{Shiyi Liu}
\IEEEauthorblockA{
Computer Science\\
Virginia Tech\\
Alexandria, VA, USA\\
\texttt{shiyiliu@vt.edu}}
}

\maketitle

\begin{abstract}
This project proposes CrisisEcho, a consumer-facing application that leverages
Retrieval-Augmented Generation (RAG) and large language model (LLM) agents to
detect, classify, and summarize hyperlocal emergencies from multi-source social
media streams in near real-time. Traditional emergency response systems rely
almost exclusively on official channels such as 911 dispatch, government sensors,
and police scanners, which frequently lag actual events by minutes or hours.
CrisisEcho addresses this gap by ingesting posts from Twitter/X, Reddit,
Bluesky, Mastodon, Nextdoor, and Telegram alongside corroborating signals from
official feeds including the National Weather Service (NWS), USGS Earthquake
Hazards, FEMA OpenFEMA, PulsePoint, and GDELT. The system stores processed
posts in MongoDB Atlas using native 2dsphere geospatial indexing, and uses a
Pinecone-hosted vector index for payload-filtered semantic retrieval. A three-step
LLM agent chain performs event clustering, severity scoring, and natural-language
alert generation. Results are delivered through a React web application featuring
a Leaflet.js interactive map with color-coded crisis hotspots, a pin detail panel,
a WebSocket real-time notification engine, and a natural-language query interface.
The backend is built in Go for high-throughput concurrent alert delivery, while
the AI pipeline runs in Python using LangChain and Sentence-Transformers. The
system is fully containerized with Docker Compose and deployed on GCP Cloud Run.
This project demonstrates how LLMs transform relational and document stores into
open-world, semantically reasoning systems capable of detecting crises from noisy,
unstructured, human-generated text.
\end{abstract}

\begin{IEEEkeywords}
emergency detection, social media mining, retrieval-augmented generation,
large language models, geospatial databases, MongoDB, real-time systems,
crisis informatics, hyperlocal awareness
\end{IEEEkeywords}

% ─────────────────────────────────────────────────────────────────────────────
\section{Introduction}
% ─────────────────────────────────────────────────────────────────────────────

\subsection{Motivation}

Traditional emergency response systems rely almost entirely on official channels,
including 911 calls, government sensors, and police scanners. These channels
frequently arrive minutes or even hours after a crisis begins. In fast-moving
events such as flash floods, active shooter situations, widespread power outages,
or natural disasters, those delays can have life-threatening consequences.

In today's connected world, ordinary people are frequently the first witnesses,
posting real-time signals on social media before any official report is filed.
Real-world examples of such posts include statements such as ``Water rising fast
on my street in Leesburg — cars already stuck!'', ``Heard multiple gunshots near
Leesburg Outlet Mall, everyone running'', ``My basement is turning into a pool
lmao — serious flooding on Maple Ave'', and ``Sounds like fireworks but way too
many — near the mall.'' These posts contain the earliest, most hyperlocal clues
about emerging crises, typically appearing before official reports are issued.

\subsection{Problem Statement}

Existing monitoring tools, both commercial and research prototypes, are
predominantly keyword-based. They scan for exact terms such as ``flood'' or
``gunshots.'' This approach has three critical failure modes. First, it misses
nuance and context, such as indirect crisis signals expressed in colloquial or
figurative language. Second, it struggles with sarcasm, slang, and regional
expressions. Third, it generates both false negatives, missing real events, and
false positives, flooding operators with noise, making such systems unreliable
for timely action.

Furthermore, most existing crisis monitoring tools are designed for professional
emergency management operators, not everyday citizens. There is no widely
accessible consumer application that synthesizes multi-source social signals,
applies LLM-based semantic reasoning, and delivers plain-language, map-based
crisis alerts to the general public.

\subsection{Significance}

CrisisEcho directly demonstrates how large language models are transforming
database systems, as highlighted in the course material. The system instantiates
four key paradigm shifts. In the dimension of Text to Semantics, semantic
understanding capabilities allow the system to discern nuances, context, and
subtleties that traditional keyword-based data management cannot. In the dimension
of Retrieval to Reasoning, multi-step LLM agent chains perform event clustering,
severity assessment, and alert generation far beyond simple lookup retrieval.
In the dimension of Vertical Domains to Multiple Domains, the pipeline is
extensible from crisis detection to positive community event detection without
schema changes. In the dimension of Closed World to Open World, the system
generalizes to evolving, unstructured real-world language without rigid schemas
or vocabulary lists.

By building CrisisEcho as a consumer-facing app, the project directly puts
earlier, more intelligent crisis awareness in the hands of everyday people,
supporting faster community coordination and personal safety decision-making.

% ─────────────────────────────────────────────────────────────────────────────
\section{Literature Review: Social Media Crisis Detection}
% ─────────────────────────────────────────────────────────────────────────────

\subsection{Social Media as a Crisis Sensor}

Social media platforms have been extensively studied as real-time data sources
for crisis detection and humanitarian response. Imran et al.~\cite{imran2016twitter}
demonstrated that Twitter contains actionable crisis-related content during
disasters and developed large-scale annotated corpora (CrisisNLP, HumAID) to
train classification models at scale. Olteanu et al.~\cite{olteanu2014crisislex}
introduced CrisisLex, a lexicon for collecting and filtering microblogged
communications during crises, establishing a foundational vocabulary approach
that CrisisEcho supersedes with semantic embeddings.

Middleton et al.~\cite{middleton2014real} demonstrated real-time crisis mapping
from social media, showing that geographically tagged social posts correlate
strongly with ground-truth disaster events. Kryvasheyeu et al.~\cite{kryvasheyeu2016rapid}
showed that social media activity can provide rapid damage assessment during
natural disasters, validating social signals as a credible early-warning source.

\subsection{Limitations of Keyword-Based Approaches}

The primary limitation of existing systems is their reliance on fixed vocabulary
matching. Such approaches consistently fail to detect crisis signals expressed
through indirect language, regional idioms, or evolving slang. Alam et
al.~\cite{alam2021humaid} established benchmarks showing that even fine-tuned
classification models trained on labeled crisis data struggle to generalize to
novel crisis types or geographic regions, motivating the open-world approach
taken by CrisisEcho.

\subsection{Machine Learning in Crisis Informatics}

Nguyen et al.~\cite{nguyen2017robust} proposed convolutional neural networks
(CNNs) for identifying valuable crisis-related information from tweets,
demonstrating improvement over keyword baselines. More recently, transformer-based
models have become the dominant architecture. Alam et al.~\cite{alam2021humaid}
showed that BERT-based classifiers trained on HumAID achieve strong F1 scores
across multiple crisis event types and classification tasks including crisis
type identification, humanitarian category labeling, and informativeness scoring.

CrisisEcho extends this line of work by replacing single-step classification
with a multi-step LLM agent reasoning chain, enabling not only detection but
also event clustering, severity scoring, and actionable alert generation in a
single unified pipeline.

% ─────────────────────────────────────────────────────────────────────────────
\section{Literature Review: LLMs, RAG, and Geospatial Database Systems}
% ─────────────────────────────────────────────────────────────────────────────

\subsection{Retrieval-Augmented Generation}

Lewis et al.~\cite{lewis2020rag} introduced Retrieval-Augmented Generation (RAG)
as a framework for combining dense retrieval with generative language models,
enabling models to ground their outputs in retrieved evidence. This approach is
central to CrisisEcho: rather than relying on LLM parametric knowledge alone,
the agent retrieves semantically similar recent posts as context before generating
a crisis assessment. This grounding reduces hallucination and anchors the LLM
output to real, timestamped evidence.

\subsection{LLM Agents and Multi-Step Reasoning}

Recent work on LLM agents, including the ReAct framework~\cite{yao2023react}
and LangChain orchestration, has demonstrated that chaining multiple LLM calls
with intermediate reasoning steps significantly outperforms single-prompt
approaches on complex analytical tasks. CrisisEcho applies a three-step agent
chain: event clustering, severity scoring with rationale, and natural-language
alert generation. This decomposition mirrors the cognitive process of a human
emergency analyst while operating at machine speed.

\subsection{Geospatial Database Systems}

MongoDB's native 2dsphere indexing supports GeoJSON geometries and enables
efficient geospatial queries including \texttt{\$near} (posts within radius of
a point), \texttt{\$geoWithin} (posts inside a polygon), and
\texttt{\$geoIntersects} (overlapping geometries). This eliminates the need for
PostGIS as a separate extension and simplifies the deployment stack. Geospatial
indexing in document databases has been shown to scale efficiently to millions
of geotagged records with sub-100ms query latency~\cite{mongodb2023geo},
meeting the real-time requirements of CrisisEcho.

\subsection{Vector Search and Hybrid Retrieval}

Karpukhin et al.~\cite{karpukhin2020dense} demonstrated that dense passage
retrieval using bi-encoder models significantly outperforms BM25 keyword
retrieval for open-domain question answering, motivating the use of dense
vector search in CrisisEcho. Pinecone's hosted vector index provides
Approximate Nearest Neighbor (ANN) search via HNSW with metadata filtering,
enabling combined vector similarity plus geographic bounding box queries in a
single retrieval step. Unlike other free-tier vector database options, Pinecone's
Starter plan provides persistent storage with no inactivity-based deletion,
making it appropriate for a semester-long course project.

% ─────────────────────────────────────────────────────────────────────────────
\section{Literature Review: Geolocation from Unstructured Text}
% ─────────────────────────────────────────────────────────────────────────────

\subsection{Text-Based Geocoding}

A fundamental challenge in social media crisis detection is that fewer than 2\%
of tweets carry GPS coordinates~\cite{middleton2014real}. The Carmen
geocoder~\cite{dredze2013carmen} addresses this by extracting location
references from post text, user profile location fields, and time zone
metadata, achieving reasonable precision on English-language social media.
spaCy's Named Entity Recognition (NER) pipeline provides a complementary
approach by identifying place names, neighborhood references, and
intersection descriptions in post text.

CrisisEcho implements a cascading geocoding strategy: GPS metadata (when
available) takes highest priority, followed by Carmen extraction, spaCy NER,
user profile location, and finally a batch inference step that uses nearby
geotagged posts as geographic priors for posts with no extractable location.

% ─────────────────────────────────────────────────────────────────────────────
\section{Dataset Description}
% ─────────────────────────────────────────────────────────────────────────────

CrisisEcho uses a multi-layered data collection strategy spanning social media
platforms, official government and emergency APIs, crowdsourced citizen platforms,
and labeled research datasets. All sources are either publicly available APIs,
academic datasets compliant with their respective terms of use, or public-facing
social media streams accessible under standard research guidelines.

\subsection{Social Media Platforms --- Live Streaming}

\textbf{Twitter/X} is accessed via Tweepy filtered stream for live ingestion
and supplemented by historical academic datasets (CrisisNLP, HumAID, Kaggle
Disaster Tweets) for training and evaluation. Twitter/X provides broad
real-time crisis signals and is the most studied platform in crisis informatics.

\textbf{Reddit} is accessed via PRAW streaming from relevant local subreddits
(r/Virginia, r/nova, r/news, r/Alexandria, r/leesburg) and supplemented by
Pushshift historical dumps. Reddit posts are longer-form and often verified by
community upvotes, making them a high-value source for geographically attributed
crisis reports.

\textbf{Bluesky} is accessed through the AT Protocol public firehose, which
imposes no rate limits and provides a fully open stream of public posts. Bluesky
is a growing platform with active local community usage and represents a
rate-limit-free alternative to Twitter/X.

\textbf{Mastodon} is accessed through instance-level API streams. As a
decentralized federated network, Mastodon provides open access to public
timelines without the restrictive API policies of centralized platforms.

\textbf{Nextdoor} public posts are collected via structured web scraping of
neighborhood feeds. Nextdoor provides extremely hyperlocal signals at the
street and block level, making it the highest-resolution geographic source
in the pipeline.

\textbf{Telegram} public emergency channels and local news channels are
ingested using the Telethon Python library. Many local emergency groups,
neighborhood channels, and breaking news accounts publish real-time incident
reports via Telegram.

\textbf{Facebook Groups} public post content is accessed via the Graph API
for public group posts. Large local community groups frequently serve as
informal crisis reporting forums.

\subsection{Official Government and Emergency APIs}

\textbf{National Weather Service (NWS)} provides a free REST API at
weather.gov delivering real-time weather alerts, severe weather warnings,
and flood advisories by geographic location. NWS alerts serve as official
corroboration signals that the LLM agent uses to increase the severity score
of social media clusters reporting weather-related events.

\textbf{USGS Earthquake Hazards API} provides free real-time earthquake data
with sub-minute latency. USGS signals cross-validate social posts reporting
shaking, building damage, or loud rumbling noises.

\textbf{FEMA OpenFEMA API} provides free access to historical disaster
declaration records and flood zone data. These are used primarily for labeling
training data and providing geographic context on historically high-risk areas.

\textbf{Waze Connected Citizens Program} makes real-time accident and road
closure data available through a public data feed, providing additional incident
corroboration for traffic-related crisis signals.

\textbf{PulsePoint} exposes live 911 dispatch incident data in participating
jurisdictions through an API reverse-engineered by the developer community.
PulsePoint data provides near-ground-truth labels for active incidents,
enabling both real-time corroboration and offline evaluation.

\subsection{Crowdsourced and Citizen Platforms}

\textbf{Ushahidi} public datasets from past global crisis deployments provide
labeled incident reports from real-world emergencies. These datasets supplement
the Twitter-centric training data with non-Twitter crisis signal examples.

\textbf{Citizen App} public incident feeds are accessed through web scraping
of the public-facing incident map. Citizen posts are directly tied to 911
dispatch data and provide strong ground-truth labels for model evaluation.

\textbf{GDELT Project} monitors global news media in near real-time and
publishes free bulk data and streaming API access. GDELT cross-references
local news signals at massive scale, providing a high-recall supplementary
source for identifying events that have been picked up by local news but
not yet widely posted on social media.

\textbf{Patch.com} RSS feeds provide hyperlocal news by locality. Patch covers
fires, accidents, power outages, and other incidents in suburban communities
at a geographic granularity well-matched to CrisisEcho's Northern Virginia focus.

\textbf{NewsAPI} and \textbf{MediaStack} free-tier REST APIs aggregate news
sources filterable by location and keyword, supporting both near-real-time
signal collection and offline ground-truth labeling of crisis events.

\subsection{Labeled Research Datasets for Training and Evaluation}

The following publicly available datasets are used for model training,
fine-tuning, and offline evaluation:

\begin{itemize}
    \item \textbf{HumAID}~\cite{alam2021humaid}: 77,000 tweets across 19 disaster
    events, annotated for crisis type, informativeness, and humanitarian category.
    \item \textbf{CrisisNLP}~\cite{imran2016twitter}: Multi-event annotated
    Twitter dataset across dozens of real-world crises with fine-grained labels.
    \item \textbf{CrisisLex}~\cite{olteanu2014crisislex}: Curated lexicon plus
    labeled tweet collections from multiple crisis events.
    \item \textbf{Kaggle Disaster Tweets}: Binary relevance-labeled tweets from
    a widely used NLP competition dataset.
    \item \textbf{CrisisMMD}: Multimodal dataset combining text and images for
    crisis classification tasks across four annotation dimensions.
    \item \textbf{FireFlood Dataset}: Labeled social media posts specific to fire
    and flood events, providing event-type-specific training examples.
    \item \textbf{ISCRAM Datasets}: Multiple labeled corpora from the Information
    Systems for Crisis Response and Management conference proceedings.
\end{itemize}

\subsection{Data Limitations}

Potential limitations of the described sources include the following. Nextdoor
and Facebook scraping may be blocked by platform anti-scraping measures,
requiring fallback to open-platform alternatives. Fewer than 2\% of Twitter
posts carry GPS coordinates, requiring text-based geocoding with inherent
uncertainty. PSM and BSM signals from connected infrastructure are not
available in the Northern Virginia region at the time of writing. Patch.com
and GDELT may have latency of several minutes before news articles are
indexed, limiting their utility as real-time signals compared to social media
streams. Despite these limitations, the combined source portfolio provides
sufficient signal density and redundancy for reliable early-warning detection
across the Northern Virginia and DC metro study area.

% ─────────────────────────────────────────────────────────────────────────────
\section{Proposed Approaches}
% ─────────────────────────────────────────────────────────────────────────────

\subsection{Multi-Source Ingestion with Kafka}

All data sources feed into a Kafka message broker that decouples ingestion from
processing. Separate Kafka topics are maintained for each source class:
\texttt{social\_raw} for all social media platforms, \texttt{official\_alerts}
for NWS, USGS, FEMA, and PulsePoint feeds, \texttt{news\_feed} for GDELT,
Patch, and NewsAPI, and \texttt{citizen\_reports} for Ushahidi and Citizen App.
Go consumers read from each topic and trigger the Python preprocessing service
via internal REST calls. For batch historical datasets (HumAID, CrisisNLP), a
Go replay worker simulates streaming at a configurable throughput rate to enable
end-to-end pipeline testing.

\subsection{Semantic Pre-filtering with DistilBERT}

Before any post enters the vector index or LLM pipeline, a lightweight
DistilBERT classifier fine-tuned on HumAID performs binary relevance
classification. This pre-filter discards approximately 70\% of incoming posts
as off-topic, reducing the volume of embeddings generated and LLM API calls
required. Only posts passing the relevance threshold proceed to geocoding and
embedding.

\subsection{Hybrid Geospatial and Vector Retrieval}

For each periodic trigger (every 60 seconds or on significant volume spike),
the retrieval module issues a hybrid query combining three concurrent lookups.
First, a Pinecone ANN search retrieves the top 50 semantically similar posts
from the last two hours filtered by geographic bounding box. Second, a MongoDB
\texttt{\$near} query retrieves all posts within a configurable radius of the
trigger centroid that may have low semantic similarity but share geographic
proximity. Third, any active official alerts from NWS, USGS, or PulsePoint
in the same area are retrieved from MongoDB and injected as high-authority
context documents. The three result sets are merged and de-duplicated before
being passed to the LLM agent.

\subsection{Three-Step LLM Agent Chain}

The core reasoning pipeline is implemented as a three-step LangChain agent
chain. In step one, the agent receives the merged retrieval results as context
and performs event clustering, identifying distinct real-world events being
described across the posts and producing a structured JSON output specifying
event type, location description, estimated time, contributing post identifiers,
and a confidence score. In step two, the agent assesses cluster severity on a
1--5 scale considering language urgency, count of independent reporters,
cross-platform corroboration, and presence of any official corroborating alert.
The agent outputs both a numeric severity score and a two-sentence rationale.
In step three, the agent generates a two-to-three sentence plain-language alert
for public consumption, ensuring no personally identifiable information appears
in the output and appending a source platform attribution list.

\subsection{Severity Thresholding and Notification}

Alerts are only published to the notification system when the cluster satisfies
three conditions simultaneously: the LLM-assigned severity score is 3 or higher
on the 1--5 scale, the cluster contains at least three independent post
contributors from at least two distinct platforms, and the cluster centroid
geocoding confidence exceeds a minimum threshold. This multi-condition filter
substantially reduces notification fatigue compared to single-threshold systems.

% ─────────────────────────────────────────────────────────────────────────────
\section{System Architecture}
% ─────────────────────────────────────────────────────────────────────────────

\subsection{Architecture Overview}

CrisisEcho is structured as a five-layer pipeline. Layer one handles
multi-source ingestion through Kafka. Layer two performs NLP preprocessing
and enrichment in Python. Layer three provides hybrid storage across MongoDB
Atlas, Pinecone, and Redis. Layer four executes the RAG-powered LLM reasoning
agent. Layer five delivers alerts and interactive visualizations through the
Go API and React frontend.

The system architecture diagram is provided below.

\begin{figure*}[h]
\centering
\includegraphics[width=\textwidth]{images/CrisisEcho Diagram.pdf}
\caption{CrisisEcho System Architecture}
\label{fig:architecture}
\end{figure*}

\subsection{Layer 1: Ingestion and Streaming}

Platform-specific ingestion workers are implemented in Python using
Tweepy (Twitter/X), PRAW (Reddit), the AT Protocol Python SDK (Bluesky),
Telethon (Telegram), and standard HTTP clients (NWS, USGS, GDELT, NewsAPI,
Patch RSS). Each worker publishes raw post objects as JSON to the appropriate
Kafka topic. Kafka consumer groups in Go handle backpressure and provide
at-least-once delivery guarantees. Dead-letter queues capture malformed or
oversized messages for offline inspection.

\subsection{Layer 2: Preprocessing and Enrichment}

A Python microservice consumes from the \texttt{social\_raw} and related
Kafka topics and applies the following sequential steps. Text cleaning
uses a spaCy pipeline with custom components for URL removal, emoji normalization
(optionally retained as sentiment signal), and whitespace normalization.
Geocoding applies Carmen as the primary extractor followed by spaCy NER for
location entity identification. If neither method yields a location, the post is
tagged as \texttt{location:unresolved} and held for batch geo-inference using
nearby geotagged posts as geographic priors. Near-duplicate detection uses
MinHash LSH via the datasketch library with a Jaccard similarity threshold of
0.85 and a five-minute deduplication window. Relevance pre-filtering applies
the DistilBERT classifier. Embedding generation uses the
\texttt{all-MiniLM-L6-v2} Sentence-Transformers model to produce 384-dimensional
dense vectors, configurable to \texttt{all-mpnet-base-v2} for higher accuracy
at the cost of throughput.

\subsection{Layer 3: Hybrid Storage}

\subsubsection{MongoDB Atlas (Primary Document and Geospatial Store)}

MongoDB Atlas free tier M0 serves as the primary document and geospatial store.
All post documents, cluster records, and alert records are stored in MongoDB.
The \texttt{posts} collection uses a 2dsphere index on the GeoJSON
\texttt{location} field to enable efficient geospatial queries. The
\texttt{clusters} collection stores both a centroid point and a convex hull
polygon for each detected crisis cluster. The \texttt{alerts} collection
maintains a full audit trail including which posts contributed to each alert
and the LLM-generated rationale.

\subsubsection{Pinecone (Vector Index)}

Pinecone's hosted free Starter plan provides a persistent vector index with
support for up to 100,000 vectors and one index at no cost. Unlike Qdrant Cloud
or Weaviate Cloud sandbox tiers, Pinecone's Starter plan does not expire
due to inactivity, ensuring that embeddings accumulated during Phase 1 remain
accessible throughout the project duration. The Pinecone index stores one vector
per processed post with payload metadata fields for post identifier, timestamp,
latitude, longitude, source platform, and crisis type, enabling combined ANN
search plus metadata filtering in a single query. The Python
\texttt{pinecone-client} library handles all vector operations; the Go API
queries Pinecone via its REST API at inference time.

\subsubsection{Redis (Hot Cache and Pub/Sub)}

Redis stores the sliding window of the most recent 30 minutes of processed
posts as a sorted set indexed by timestamp. LLM-generated cluster summaries
are cached with a five-minute TTL to prevent redundant API calls for
geographically overlapping retrieval queries. The Redis Pub/Sub channel
\texttt{alerts:live} carries new alert events from the Python agent to the Go
API server, which relays them to connected React clients over WebSocket.

\subsection{Layer 4: RAG Pipeline and LLM Agent}

The RAG pipeline is implemented in Python using LangChain for agent
orchestration and LlamaIndex for document loading and context window
management. The LLM backend is Claude claude-haiku-4-5 (primary, optimized for
speed and cost) with Ollama-hosted Llama3 as a configurable free fallback
for offline or budget-constrained operation. The agent is stateless; all
relevant context is included in each API call as retrieved documents.
New cluster documents are written back to MongoDB upon successful completion
of all three agent steps, and a Redis Pub/Sub message is published to trigger
downstream alert delivery.

\subsection{Layer 5: Go API and React Frontend}

The Go backend uses the Fiber HTTP framework and exposes the following REST
endpoints. \texttt{GET /api/hotspots} accepts latitude, longitude, and radius
parameters and returns all active crisis clusters within the specified area.
\texttt{GET /api/pin} returns detailed cluster information for a user-dropped
map pin including the LLM summary, contributing posts, timeline, and source
breakdown. \texttt{POST /api/subscribe} stores user notification preferences
(location, crisis type filters, radius). \texttt{GET /api/query} accepts a
natural-language text query, routes it to the LLM agent with a specialized
digest prompt, and returns a structured response. \texttt{GET /ws/alerts}
is a WebSocket endpoint that streams new alert events to connected clients
using Redis Pub/Sub as the backend message bus.

The React frontend uses TypeScript and react-leaflet for the interactive map.
Crisis hotspots are rendered as circle markers sized proportionally to severity
score and colored by crisis type (red for life-safety, orange for
infrastructure, yellow for traffic, green for positive events in the stretch
goal). A slide-up pin detail drawer shows the LLM summary, a Recharts timeline
of post volume over time, anonymized post excerpts with source labels, and
cross-platform corroboration indicators. A WebSocket hook subscribes to
\texttt{/ws/alerts} and renders toast notifications for new high-severity events.
A natural-language query bar allows users to type questions such as ``What is
happening near Dulles right now?'' and receive an LLM-generated digest.

\subsection{Infrastructure and Deployment}

All services are defined in a single Docker Compose file for local development
and end-to-end testing. Production deployment targets GCP Cloud Run, which
provides serverless container hosting with a generous free tier. MongoDB Atlas
and Pinecone are externally hosted and require no infrastructure management.
The Go API and Python preprocessing service are containerized separately,
enabling independent scaling. The React frontend is built to a static bundle
and hosted on Firebase Hosting or GCP Cloud Storage.

% ─────────────────────────────────────────────────────────────────────────────
\section{Database Schema Design}
% ─────────────────────────────────────────────────────────────────────────────

\subsection{MongoDB: posts Collection}

Each document in the \texttt{posts} collection contains the following fields.
The \texttt{\_id} field is a MongoDB ObjectId. The \texttt{post\_id} field is
a string in the format \texttt{platform:native\_id}
(e.g., \texttt{twitter:1234567890}). The \texttt{source} field identifies
the platform. The \texttt{raw\_text} and \texttt{cleaned\_text} fields store
the original and preprocessed content respectively. The \texttt{timestamp}
field is an ISODate. The \texttt{location} field is a GeoJSON Point object
\texttt{\{type: "Point", coordinates: [lng, lat]\}} covered by a 2dsphere index.
The \texttt{location\_confidence} is a float from 0 to 1 indicating geocoding
certainty. The \texttt{location\_source} field identifies whether the location
was derived from GPS, Carmen, NER, user profile, or inference. The
\texttt{embedding\_id} field references the corresponding Pinecone vector.
The \texttt{cluster\_id} is a nullable ObjectId referencing the
\texttt{clusters} collection. The \texttt{crisis\_type} field is a nullable
categorical string from the set: flood, shooting, outage, fire, accident,
earthquake, or null. The \texttt{severity\_score} is an integer from 1 to 5.
The \texttt{is\_relevant} field is a boolean output of the DistilBERT pre-filter.
The \texttt{metadata} subdocument stores a hashed username and platform-specific
metadata.

\subsection{MongoDB: clusters Collection}

Each document in the \texttt{clusters} collection contains the following fields.
The \texttt{centroid} field is a GeoJSON Point covered by a 2dsphere index. The
\texttt{affected\_area} field is a GeoJSON Polygon representing the convex hull
of all contributing post locations. The \texttt{crisis\_type}, \texttt{severity},
and \texttt{severity\_rationale} fields store the LLM agent outputs from step
two of the reasoning chain. The \texttt{summary} field stores the step-three
alert text. The \texttt{post\_count} and \texttt{post\_ids} fields track cluster
membership. The \texttt{sources} array lists all contributing platforms. The
\texttt{official\_corroboration} boolean indicates whether any NWS, USGS, or
PulsePoint signal was present in the retrieval context. The \texttt{start\_time},
\texttt{last\_updated}, and \texttt{status} fields (active, resolved, or stale)
track cluster lifecycle.

\subsection{Pinecone: post\_embeddings Index}

The Pinecone index stores 384-dimensional float32 vectors generated by
\texttt{all-MiniLM-L6-v2}. Each vector is associated with a metadata payload
containing: \texttt{post\_id}, \texttt{timestamp} (Unix epoch), \texttt{lat},
\texttt{lng}, \texttt{source}, \texttt{crisis\_type}, and
\texttt{severity\_score}. All metadata fields are indexed for filtering,
enabling combined ANN plus metadata queries such as ``find the top 50 semantically
similar posts from the last two hours within a 50km radius.''

% ─────────────────────────────────────────────────────────────────────────────
\section{RAG Pipeline: Detailed Specification}
% ─────────────────────────────────────────────────────────────────────────────

\subsection{Retrieval Strategy}

For each pipeline trigger, the retrieval module issues three concurrent queries.
The Pinecone query uses HNSW ANN search to retrieve the top 50 vectors with
\texttt{timestamp} greater than current time minus 7200 seconds and
\texttt{lat}/\texttt{lng} within the trigger bounding box. The MongoDB query
uses \texttt{\$near} with a configurable radius (default 50km) to retrieve
posts with low semantic similarity but high geographic proximity. Official
signals from the \texttt{official\_alerts} Kafka topic that have been persisted
to a dedicated MongoDB collection are retrieved using a spatial intersection
query. Results are merged, de-duplicated on \texttt{post\_id}, and ranked by
a composite score combining semantic similarity, recency, and source authority.

\subsection{LLM Prompt Templates}

The step-one clustering prompt instructs the model: ``Given these N social
media posts from the past two hours near [location], identify distinct
real-world events being described. For each event output: event\_type,
location\_description, estimated\_time, contributing\_post\_ids, and
confidence\_score as JSON.''

The step-two severity prompt instructs the model: ``Given this cluster of N posts
describing [event\_type] near [location], rate severity 1--5 where 1 is
unconfirmed minor incident and 5 is confirmed mass casualty or critical
infrastructure failure. Consider language urgency, number of independent
reporters, and any official corroboration. Output: severity\_score (integer)
and rationale (two sentences) as JSON.''

The step-three alert prompt instructs the model: ``Write a two-to-three sentence
alert for general public consumption about this severity-[X] [event\_type]
event near [location]. Be factual, calm, and actionable. Do not include
usernames. Append: Sources: [platform list].''

% ─────────────────────────────────────────────────────────────────────────────
\section{Evaluation Methodology}
% ─────────────────────────────────────────────────────────────────────────────

\subsection{Baseline Comparison}

CrisisEcho is evaluated against a keyword-based baseline that applies the
CrisisLex lexicon to flag posts matching any crisis-related term. The same
post corpus, label set, and evaluation split are used for both systems,
enabling a direct comparison of precision, recall, and F1.

\subsection{Evaluation Metrics}

The primary evaluation metrics are as follows. Precision at K measures the
fraction of the K fired alerts that correspond to a verified real-world event
in the ground truth, with a target of greater than 0.75. Recall measures the
fraction of ground-truth events that CrisisEcho detected with at least one
alert, with a target of greater than 0.70. End-to-end latency measures the
time from the first social post about an event to the alert appearing on the
map, with a target of under 90 seconds. The false positive rate measures the
fraction of fired alerts with no corresponding real event, with a target below
0.20. Geocoding accuracy measures the fraction of posts geocoded to the correct
city or neighborhood, with a target above 0.80. The F1 improvement over the
keyword baseline is expected to exceed 15 percentage points based on comparable
systems in the literature.

\subsection{Ground Truth Construction}

Ground truth labels are constructed from two sources. Offline evaluation uses
the HumAID and CrisisNLP test splits with their existing event labels.
Online evaluation during the project period uses PulsePoint dispatch records
and NWS active alerts as independent ground truth for Northern Virginia events,
enabling evaluation of the full live pipeline.

% ─────────────────────────────────────────────────────────────────────────────
\section{Challenges and Mitigations}
% ─────────────────────────────────────────────────────────────────────────────

\textbf{Noisy Unstructured Text.} Social media contains typos, sarcasm, memes,
and off-topic posts at high volume. The DistilBERT pre-filter removes approximately
70\% of noise before LLM processing, and MinHash LSH deduplication removes
reposts and retweets.

\textbf{Sparse Geolocation.} The Carmen and spaCy NER cascade with profile location
fallback addresses most cases. Posts with no extractable location are held for
batch geo-inference using nearby geotagged posts as geographic priors rather
than being discarded.

\textbf{LLM API Cost.} The Claude Haiku model is the lowest-cost Anthropic API
tier. Posts are batched before each API call. Redis caches cluster summaries for
five minutes. Llama3 via Ollama provides a free offline fallback.

\textbf{Pinecone Free Tier Limits.} The Starter plan provides 100,000 vectors.
Embeddings older than 72 hours are deleted via a scheduled cleanup worker,
maintaining headroom within the free tier throughout the project.

\textbf{Real-Time Latency.} Kafka batching, Go concurrency, and Redis caching
keep end-to-end latency under the 90-second target. The DistilBERT pre-filter
reduces the volume of posts requiring embedding and LLM processing by approximately
70\%.

\textbf{False Positives and Notification Fatigue.} The three-condition alert
threshold (severity 3 or higher, at least three independent contributors from
at least two platforms, minimum geocoding confidence) substantially reduces
false positive alerts compared to single-threshold systems.

\textbf{Privacy and Ethics.} All usernames are hashed at ingestion using SHA-256
before storage. No personally identifiable information appears in LLM prompts
or alert outputs. Notifications are strictly opt-in. The system is designed for
situational awareness and safety, not surveillance.

\textbf{Platform Access Restrictions.} Nextdoor and Facebook scraping may be
blocked. Both are classified as supplementary sources. Bluesky, Mastodon, and
Reddit provide equivalent open-API coverage as fallback alternatives.

% ─────────────────────────────────────────────────────────────────────────────
\section{Technology Stack Summary}
% ─────────────────────────────────────────────────────────────────────────────

Table~\ref{tab:techstack} summarizes the complete technology stack with
rationale for each selection.

\begin{table}[h]
\caption{CrisisEcho Technology Stack}
\label{tab:techstack}
\renewcommand{\arraystretch}{1.2}
\begin{tabular}{p{1.4cm} p{2.2cm} p{4.0cm}}
\toprule
\textbf{Component} & \textbf{Technology} & \textbf{Rationale} \\
\midrule
Backend API & Go (Fiber) & High-throughput concurrent alert delivery; goroutines for WebSocket management \\
AI Pipeline & Python 3.11 + LangChain & Rich NLP/LLM ecosystem; LangChain for multi-step agent orchestration \\
Frontend & React 18 + TypeScript + Leaflet.js & Component-based UI; react-leaflet for interactive geospatial map \\
Primary DB & MongoDB Atlas (Free M0) & Document store with native 2dsphere geospatial indexing; no PostGIS needed \\
Vector DB & Pinecone (Free Starter) & Persistent hosted ANN search; payload-filtered vector queries; no inactivity expiry \\
Cache / Pub-Sub & Redis 7 & Sliding post window; LLM result cache; WebSocket alert relay \\
Message Queue & Apache Kafka & Decouples ingestion from processing; simulates real-time streaming \\
NLP & spaCy + Sentence-Transformers + datasketch & Cleaning, NER, embeddings, MinHash LSH deduplication \\
Geocoding & Carmen + spaCy NER & Cascading text-based location extraction from unstructured posts \\
LLM Backend & Claude claude-haiku-4-5 / Ollama Llama3 & Haiku for speed/cost; Llama3 as free offline fallback \\
Containers & Docker + Docker Compose & Single-command local dev; reproducible environment \\
Deployment & GCP Cloud Run & Serverless, pay-per-use, free tier covers prototype traffic \\
Visualization & Recharts + Leaflet.js & Timeline charts and interactive geospatial map in React \\
\bottomrule
\end{tabular}
\end{table}

% ─────────────────────────────────────────────────────────────────────────────
\begin{thebibliography}{00}

\bibitem{imran2016twitter}
M. Imran, P. Mitra, and C. Castillo, ``Twitter as a lifeline: Human-annotated
Twitter corpora for NLP of crisis-related messages,'' in \textit{Proc. 10th Int.
Conf. Language Resources and Evaluation (LREC)}, 2016, pp. 1638--1643.

\bibitem{olteanu2014crisislex}
A. Olteanu, C. Castillo, F. Diaz, and S. Vieweg, ``CrisisLex: A lexicon for
collecting and filtering microblogged communications in crises,'' in \textit{Proc.
8th Int. AAAI Conf. Weblogs and Social Media (ICWSM)}, 2014, pp. 376--385.

\bibitem{alam2021humaid}
F. Alam, F. Sajjad, H. Dalvi, N. Durrani, K. Darwish, and P. Nakov, ``HumAID:
Human-annotated disaster incidents data from Twitter with deep learning
benchmarks,'' in \textit{Proc. 15th Int. AAAI Conf. Weblogs and Social Media
(ICWSM)}, 2021.

\bibitem{middleton2014real}
S. E. Middleton, L. Middleton, and S. Modafferi, ``Real-time crisis mapping of
natural disasters using social media,'' \textit{IEEE Intelligent Systems},
vol. 29, no. 2, pp. 9--17, 2014.

\bibitem{kryvasheyeu2016rapid}
Y. Kryvasheyeu, H. Chen, N. Obradovich, E. Moro, P. Van Hentenryck,
J. Fowler, and M. Cebrian, ``Rapid assessment of disaster damage using social
media activity,'' \textit{Science Advances}, vol. 2, no. 3, p. e1500779, 2016.

\bibitem{lewis2020rag}
P. Lewis, E. Perez, A. Piktus, F. Petroni, V. Karpukhin, N. Goyal, H. Kuebler,
M. Lewis, W. Yih, T. Rocktaschel, and S. Riedel, ``Retrieval-augmented generation
for knowledge-intensive NLP tasks,'' in \textit{Advances in Neural Information
Processing Systems (NeurIPS)}, vol. 33, 2020, pp. 9459--9474.

\bibitem{yao2023react}
S. Yao, J. Zhao, D. Yu, N. Du, I. Shafran, K. Narasimhan, and Y. Cao,
``ReAct: Synergizing reasoning and acting in language models,'' in \textit{Proc.
11th Int. Conf. Learning Representations (ICLR)}, 2023.

\bibitem{karpukhin2020dense}
V. Karpukhin, B. Oguz, S. Min, P. Lewis, L. Wu, S. Edunov, D. Chen, and
W. Yih, ``Dense passage retrieval for open-domain question answering,'' in
\textit{Proc. 2020 Conf. Empirical Methods in Natural Language Processing
(EMNLP)}, 2020, pp. 6769--6781.

\bibitem{nguyen2017robust}
D. T. Nguyen, K. A. Al Mannai, S. Joty, H. Sajjad, M. Imran, and P. Mitra,
``Robust classification of crisis-related data on social networks using
convolutional neural networks,'' in \textit{Proc. 11th Int. AAAI Conf. Weblogs
and Social Media (ICWSM)}, 2017, pp. 632--635.

\bibitem{dredze2013carmen}
M. Dredze, M. Paul, S. Bergsma, and H. Tran, ``Carmen: A Twitter geolocation
system with applications to public health,'' in \textit{AAAI Workshop on
Expanding the Boundaries of Health Informatics Using AI}, 2013.

\bibitem{mongodb2023geo}
MongoDB Inc., ``Geospatial Queries,'' MongoDB Documentation, 2023. [Online].
Available: https://www.mongodb.com/docs/manual/geospatial-queries/

\bibitem{schultz2025}
G. Schultz, ``Prioritizing safety funding using severity weighted risk scores,''
\textit{Accident Analysis \& Prevention}, 2025.

\bibitem{hossain2019real}
M. Hossain, M. Abdel-Aty, M. A. Quddus, Y. Muromachi, and S. N. Sadeek,
``Real-time crash prediction models: State-of-the-art, design pathways and
ubiquitous requirements,'' \textit{Accident Analysis \& Prevention}, vol. 124,
pp. 66--84, 2019.

\end{thebibliography}

% ─────────────────────────────────────────────────────────────────────────────
\appendix
% ─────────────────────────────────────────────────────────────────────────────

\section{Task Assignment for Each Member}

\begin{table}[h]
\caption{Task Assignment Table}
\label{tab:tasks}
\renewcommand{\arraystretch}{1.2}
\begin{tabular}{p{5.0cm} p{2.5cm}}
\toprule
\textbf{Task} & \textbf{Assigned To} \\
\midrule
Define objectives, KPIs, and evaluation metrics & C. Williams \\
Draft system architecture and data flow diagram & C. Williams \\
Set repository structure, naming conventions, meeting cadence & C. Williams \\
Write Ethics, Privacy, and Anonymization section & C. Williams \\
Final report writing and figure integration & C. Williams \\
Ingest all data sources; standardize timestamps and identifiers & A. Sivakumar \\
Basic data QC (missing rates, ranges, anomalies) & Shiyi Liu \\
Data dictionary (columns, dtypes, units, sources) & Shiyi Liu \\
Implement spaCy NLP pipeline and Carmen geocoding & A. Sivakumar \\
Implement MinHash LSH deduplication & A. Sivakumar \\
Implement DistilBERT relevance pre-filter fine-tuning & A. Sivakumar \\
Define MongoDB schema and 2dsphere indexes & S. Zih \\
Implement Pinecone vector index and embedding pipeline & S. Zih \\
Build LangChain three-step LLM agent chain & S. Zih \\
Create Kafka topics and Go consumer workers & A. Sivakumar \\
Implement hybrid retrieval (Pinecone ANN + MongoDB \$near) & S. Zih \\
Integrate official signal injection (NWS, USGS, PulsePoint) & S. Zih \\
Build Go Fiber REST API and WebSocket server & S. Zih \\
Build React frontend map, hotspot markers, pin detail drawer & Ki Hong Park \\
Implement Redis Pub/Sub for real-time alert relay & S. Zih \\
Implement user notification settings and subscription logic & Ki Hong Park \\
Implement natural-language query bar (frontend + API) & Ki Hong Park \\
Run evaluation suite (Precision@K, Recall, latency, geo accuracy) & Shiyi Liu \\
Sensitivity analysis of severity threshold parameters & Shiyi Liu \\
Ablations: compare RAG vs. keyword baseline on CrisisNLP split & Shiyi Liu \\
Dashboard v2 (legends, tooltips, export PNG/CSV) & Ki Hong Park \\
Docker Compose integration test (all services, single command) & A. Sivakumar \\
GCP Cloud Run deployment and Firebase frontend hosting & C. Williams \\
Repository cleanup and README reproduction guide & A. Sivakumar \\
Demo preparation: crisis scenario replay and live walkthrough & Ki Hong Park \\
Final proof (numbering, captions, references) & Shiyi Liu \\
Package submission (PDF + appendices + code link) & Shiyi Liu \\
\bottomrule
\end{tabular}
\end{table}

\section{Project Schedule}

\begin{table}[h]
\caption{Project Schedule (15 Weeks)}
\label{tab:schedule}
\renewcommand{\arraystretch}{1.3}
\begin{tabular}{p{0.8cm} p{6.6cm}}
\toprule
\textbf{Week} & \textbf{Task} \\
\midrule
1 & Literature Review and Requirements Finalization: Review social media crisis
detection, RAG, and geospatial database literature; finalize project scope,
datasets, and technical stack (Go, Python, React, MongoDB, Pinecone, Kafka). \\
2 & Infrastructure Setup: Provision MongoDB Atlas M0 free cluster, Pinecone
Starter index, Redis Cloud free tier, and Kafka via Docker Compose; define all
collection schemas and indexes. \\
3 & Data Ingestion and Preprocessing: Implement all platform ingestion workers
(Tweepy, PRAW, AT Protocol, Telethon, NWS, USGS, GDELT, Patch RSS, NewsAPI);
implement spaCy NLP pipeline, Carmen geocoding, and MinHash deduplication;
ingest 10,000 labeled posts from HumAID and CrisisNLP as baseline corpus. \\
4 & Embedding Pipeline and Pre-filter: Implement Sentence-Transformers embedding
generation; fine-tune DistilBERT relevance classifier on HumAID; write all
processed posts and embeddings to MongoDB and Pinecone. \\
5 & RAG Retrieval Module: Implement hybrid retrieval combining Pinecone ANN
search with MongoDB \$near geospatial queries; integrate official signal
injection for NWS, USGS, and PulsePoint alerts. \\
6 & LLM Agent Chain: Build LangChain three-step agent chain (cluster, severity,
alert); integrate Claude Haiku API; configure Llama3 Ollama fallback; test
end-to-end on replayed HumAID subset. \\
7 & Cluster Persistence and Redis Integration: Write LLM agent outputs to MongoDB
clusters collection; implement Redis caching for cluster summaries;
implement Pub/Sub publish for new alert events. \\
8 & Go API Development: Build Fiber REST API with all endpoints (/api/hotspots,
/api/pin, /api/query, /api/subscribe) and WebSocket /ws/alerts server backed
by Redis Pub/Sub. \\
9 & React Frontend Development: Build Leaflet map with hotspot markers, pin
detail drawer with Recharts timeline, notification settings modal, and
natural-language query bar with WebSocket alert toasts. \\
10 & Evaluation Suite: Run offline evaluation on HumAID and CrisisNLP test splits;
compute Precision@K, Recall, F1, geocoding accuracy, and end-to-end latency;
compare RAG pipeline against CrisisLex keyword baseline. \\
11 & Stretch Features: Implement positive event detection with green markers;
begin multi-language support using multilingual-e5 embeddings for
Spanish-language posts in the Northern Virginia area. \\
12 & System Testing and Optimization: Full Docker Compose integration test; address
latency bottlenecks; tune severity thresholds; harden edge cases in
geocoding pipeline and data ingestion workers. \\
13 & Refinement and Documentation: Adjust LLM prompt templates based on evaluation
results; prepare technical documentation, API reference, and data dictionary. \\
14 & Demo Preparation: Prepare demo script simulating a flooding crisis from
replayed HumAID data; verify full pipeline fires an alert within 90 seconds;
prepare presentation slides and architecture diagram. \\
15 & Finalization: Final report writing; repository cleanup with reproduction README;
package submission (PDF, appendices, code link, data dictionary). \\
\bottomrule
\end{tabular}
\end{table}

\section{Acknowledgment}

This proposal used AI assistance for content generation and will use AI
assistance in writing code throughout the project. All architectural decisions,
specification choices, and evaluation designs were formulated by the project
team.

\end{document}